\section{Markov DGP Results}\label{sec:v10-markov-results}

The manifesto results test the framework on an expert benchmark. We
also test it on a controlled DGP where the oracle is constructible
and C1, C2, and C3 can be checked at every node without an external
expert. The Markov changepoint ops-count family generates token
streams from a finite-state
machine whose hidden state evolves by sparse changepoints; the oracle
\(f^\ast\) reads off the count of state transitions of a designated
ops type. C1 sufficiency is checkable at the leaf level, C2
idempotence is checkable on any produced state, and C3 merge
consistency reduces to an algebra on transition counts.

We use two benchmark families:
\textbf{recoverable\_v4} is the easier benchmark, where local segment
counts already encode the oracle, and
\textbf{structural\_core\_v1::r12\_seg10to12} is the harder benchmark,
where the oracle depends on cross-segment interactions and a flat
read of leaf segments cannot reconstruct it. The flat single-leaf
geometry on the harder benchmark is the regime where local laws are
expected to matter most. The tree learner is the same unified \(g\)
architecture used in the manifesto experiments. The flat baseline is
the official Fourier neural operator FNO family
\citep{Kovachki23NeuralOperator,Li21FNO}, which sees the document as one
continuous signal and has no leaf or merge supervision channel.

\subsection{Tree--FNO Parity at Full Root Supervision}

Table~\ref{tab:v10-markov-parity} compares tree-neural under the full
local-law package with the matched FNO baseline at full root
supervision (every training document carries a root oracle label).

\begin{table}[t]
\centering
\small
\setlength{\tabcolsep}{6pt}
\begin{tabular}{@{}lrrrr@{}}
\toprule
Benchmark & Geometry (leaves/doc) & Tree root MAE & FNO root MAE & \(\Delta\) \\
\midrule
recoverable\_v4 & 8 (multi-leaf) & 0.0000 & 0.0000 & \(\phantom{-}0.000\) \\
recoverable\_v4 & 1 (flat parity) & 0.0026 & 0.0026 & \(\phantom{-}0.000\) \\
structural\_core::r12 & 16 (multi-leaf) & 0.182 & 0.440 & \(-0.258\) \\
structural\_core::r12 & 4 (multi-leaf) & 0.083 & 0.096 & \(-0.014\) \\
structural\_core::r12 & 1 (flat parity) & 0.678 & 0.642 & \(\phantom{-}0.036\) \\
\bottomrule
\end{tabular}
\caption{\textbf{Tree--FNO parity at full root supervision.} Tree
runs use the local-law package (C1+C2+C3 with depth-equal weights);
FNO runs use the canonical \texttt{official\_fno\_sumlen}
comparator. \(\Delta=\text{Tree}-\text{FNO}\), so negative values
favor the tree. Source:
\texttt{outputs/markov\_v3\_publication\_fullval\_20260412\_0142}.}
\label{tab:v10-markov-parity}
\end{table}

Three patterns:

\textbf{Recoverable benchmark.} Both architectures hit zero MAE.
Local segment counts carry the oracle, and either family recovers
it. This is the easy regime; it confirms that the tree's extra
machinery does not pay a cost when the flat baseline is already
adequate.

\textbf{Structural benchmark, multi-leaf.} The tree cuts FNO root
MAE by 0.26 at 16 leaves/doc (0.182 vs 0.440). This is the regime
where C3 merge consistency carries information that a flat
representation cannot. The advantage shrinks to 0.014 at 4 leaves
because some of the relevant interactions now sit within a single
leaf and the FNO sees them.

\textbf{Structural benchmark, flat parity.} At a single leaf per
document the tree degenerates to the FNO geometry, and the two
families are within 0.04 MAE. This is the negative control: the
multi-leaf advantage above is not a generic architectural
preference, it is the local-law channel doing work where the DGP
needs it.

\subsection{Law-Package Ablation}

Table~\ref{tab:v10-markov-packages} compares six law packages on
the structural benchmark. Each package fixes a different
combination of supervision channels.

\begin{table}[t]
\centering
\small
\setlength{\tabcolsep}{5pt}
\begin{tabular}{@{}lp{0.40\linewidth}rr@{}}
\toprule
Package & Active channels & Multi-leaf MAE & Flat-parity MAE \\
\midrule
\texttt{leaf\_only} & C1 only (leaf--text supervision) & 0.111 & --- \\
\texttt{depth\_equal} & root + C1+C2+C3 with equal weights at every depth & 0.187 & --- \\
\texttt{local\_law} & root + C1+C2+C3 with depth discount & 0.182 & --- \\
\texttt{r100\_superset\_local10} & root + 10\% sampled local laws & 0.168 & --- \\
\texttt{root\_budget\_multileaf} & root supervision only, multi-leaf geometry & 0.102 & --- \\
\texttt{oneleaf\_root\_budget} & root supervision only, single-leaf geometry & --- & 0.678 \\
\bottomrule
\end{tabular}
\caption{\textbf{Law-package ablation on the structural benchmark.}
Multi-leaf cells use the best non-flat geometry per package;
flat-parity cells use a single leaf per document. The MAE values
are the headline test root MAE on the
\texttt{structural\_core::r12\_seg10to12} scope. Source:
\texttt{outputs/markov\_v3\_publication\_fullval\_20260412\_0142}.}
\label{tab:v10-markov-packages}
\end{table}

\texttt{root\_budget\_multileaf} performs surprisingly well at
0.102 MAE: the tree geometry alone, even without local-law
supervision, already cuts the structural error relative to the flat
\texttt{oneleaf\_root\_budget} baseline at 0.678. Adding local laws
on top of the same multi-leaf geometry holds the gain in the same
range and adds the auditable channel that the certificate uses.
\texttt{leaf\_only} (C1 alone) at 0.111 shows that leaf sufficiency
is the load-bearing law on this DGP; the merge and idempotence
channels add headroom but are not required for first-order recovery.

\subsection{Leaf-Size Sweep}

Figure~\ref{fig:v10-markov-leaf-size} reports the supervision-recovery
heatmap across leaf-size geometries (8 to 64 tokens per leaf) and
root-supervision shares.

\begin{figure}[t]
\centering
\begin{minipage}[t]{0.48\linewidth}
\centering
\includegraphics[width=\linewidth]{markov_leaf_size_fixed_recoverable.pdf}
\\\textbf{recoverable\_v4}
\end{minipage}\hfill
\begin{minipage}[t]{0.48\linewidth}
\centering
\includegraphics[width=\linewidth]{markov_leaf_size_fixed_structural.pdf}
\\\textbf{structural\_core::r12}
\end{minipage}
\caption{\textbf{Leaf-size sweep.} Each cell reports root MAE for a
given leaf token budget (rows) and root-supervision share
(columns). Recoverable saturates near zero across most cells;
structural rewards smaller leaves and higher root coverage, with
the local-law package lifting the bottom row. Source:
\texttt{outputs/markov\_v3\_publication\_leaf\_size\_fixed\_docs\_20260414\_200309}.}
\label{fig:v10-markov-leaf-size}
\end{figure}

\subsection{Cross-DGP Takeaway}

The Markov DGP and the manifesto data are very different objects.
Markov sentences are tokens from a state machine; manifestos are
party platforms. Markov oracles are exactly computable; manifesto
oracles are noisy expert means. Yet the same framework recovers the
same shape of result: when the DGP requires interactions across
leaves, multi-leaf trees with local laws beat flat baselines; when
local information is enough, both architectures match; when the
geometry collapses to a single leaf, the architectures are
indistinguishable. The local-law channel is the part that
generalizes across these settings, and the certificate of
Section~\ref{sec:v8-audit-certificate} is the part that quantifies
it on either DGP.
